\section{Separator information is the obstruction}\label{sec:sharpness}

The upper bounds of \cref{sec:kernels,sec:ordinary} hold however the local
functionals are chosen. This section shows that, for the full preordering,
the inverse-square rate cannot be improved in general, even for a quadratic
objective in three variables, and that the obstruction does not come from
local positivity. We exhibit a fixed problem with two bags of size two whose
sparse gap is of exact order $r^{-2}$ (\cref{thm:quad-sharp}). The lower
bound is witnessed by \emph{actual} local probability measures whose
separator moments agree through degree $2r$. It therefore applies to every
relaxation whose feasible set contains these integration functionals and
whose only coupling between bags is agreement of these separator moments.
The dense preordering on the same problem is
exact at order two. For the first two orders we compute the sparse values
exactly (\cref{prop:exact-orders}).

Minimizing each bag over its private coordinate leaves the two separator
functions $-\varphi(y)$ and $\varphi(y)$, where
$\varphi(y)=\max\{y,0\}^2$. Two separator laws with equal
moments through degree $n$ cannot tell $\varphi$ apart from its best
polynomial approximation of degree $n$, and the best approximation error of
$\varphi$ decays exactly like $n^{-2}$. The example shows that finite
separator information can limit the rate even with exact local measures; it
does not determine the rate for every sparse problem.

\subsection{Moment matching and uniform approximation}
\label{sec:shp-matching}

For a compact interval $I$, a function $\psi\in C(I)$ and $n\ge0$, let
$E_n(\psi;I)=\min\{\norm{\psi-q}_{\infty,I}:\ q\in\R[y],\ \deg q\le n\}$,
and write $E_n(\psi)=E_n(\psi;[-1,1])$.

\begin{lemma}[Moment matching]\label{lem:moment-matching-general}
Let $I$ be a compact interval, $\psi\in C(I)$, and $n\ge0$. Then
\[
  \max\Bigl\{\int\psi\,d\mu_1-\int\psi\,d\mu_2\Bigr\}=2E_n(\psi;I),
\]
where the maximum is over pairs of Borel probability measures $\mu_1,\mu_2$
on $I$ with $\int y^k\,d\mu_1=\int y^k\,d\mu_2$ for $k=0,\dots,n$, and it
is attained.
\end{lemma}

\begin{proof}
For any such pair and any polynomial $q$ of degree at most $n$,
$\int\psi\,d\mu_1-\int\psi\,d\mu_2=\int(\psi-q)\,d\mu_1-\int(\psi-q)\,d\mu_2
\le2\norm{\psi-q}_\infty$. Minimizing over $q$ gives the upper bound.

If $E_n(\psi;I)=0$, take $\mu_1=\mu_2$. Otherwise, the polynomials of
degree at most $n$ form a finite-dimensional, hence closed, subspace $P_n$
of $C(I)$ at distance $E_n(\psi;I)$ from $\psi$. By the Hahn--Banach
theorem there is a linear functional $\ell$ on $C(I)$ of norm one that
vanishes on $P_n$ and satisfies $\ell(\psi)=E_n(\psi;I)$. By the Riesz
representation theorem~\cite{rudin1987-real-complex-analysis},
$\ell(g)=\int g\,d\lambda$ for a finite signed Borel measure $\lambda$ of
total variation one. Since $\ell(1)=0$, the Jordan parts $\lambda_\pm$ have
equal mass $\frac12$. The probability measures $\mu_1=2\lambda_+$ and
$\mu_2=2\lambda_-$ satisfy $\int y^k\,d\mu_1-\int y^k\,d\mu_2=2\ell(y^k)=0$
for $k\le n$, and $\int\psi\,d\mu_1-\int\psi\,d\mu_2=2\ell(\psi)=
2E_n(\psi;I)$.
\end{proof}

This identity is known; it is stated explicitly, for example, by Han, Jiao
and Weissman~\cite{han2018-local-moment-matching}, Lemma~25, and we include
it for completeness. The lower bounds below do not need it: they use the
explicit witness of \cref{lem:fejer-lower}.

\subsection{The example}\label{sec:shp-example}

Consider
\begin{equation}\label{eq:shp-problem}
  \fstar=\min\bigl\{f(x,y,z):\ x,z\in[0,1],\ y\in[-1,1]\bigr\},\qquad
  f=\underbrace{x^2-2xy}_{f_1(x,y)}
   +\underbrace{y^2+z^2+2yz}_{f_2(y,z)},
\end{equation}
with bags $B_1=\{x,y\}$ and $B_2=\{y,z\}$ and separator $\{y\}$. The
affine change $x=(1+\hat x)/2$, $z=(1+\hat z)/2$ maps the problem to
$[-1,1]^3$, keeps both bags and all total degrees, and maps the generators
$1-\hat x^2$, $1-\hat z^2$ to $4g_x$, $4g_z$, where
\[
  g_x=x(1-x),\qquad g_y=1-y^2,\qquad g_z=z(1-z).
\]
Thus the sparse hierarchies of \cref{def:sparse-hierarchy} for the
transformed problem are, in the original coordinates, the hierarchies whose
cones are generated by $g_x,g_y$ on $B_1$ and by $g_y,g_z$ on $B_2$, with
every certificate summand of total degree at most $2r$. We write
$\rhomod_r$ and $\rhopre_r$ for their moment values on this instance.

Since $f=(y-x+z)^2+2xz$, we have $f\ge0$ on the box and $\fstar=f(0,0,0)=0$.
Let $\varphi(y)=\max\{y,0\}^2$. For every $y\in[-1,1]$,
\begin{equation}\label{eq:shp-fibers}
  \min_{x\in[0,1]}f_1(x,y)=-\varphi(y),\qquad
  \min_{z\in[0,1]}f_2(y,z)=\varphi(y),
\end{equation}
attained at $x=\max\{y,0\}$ and $z=\max\{-y,0\}$. Indeed
$f_1=(x-y)^2-y^2$ and $f_2=(y+z)^2$, and these choices lie in $[0,1]$.

For $n\ge0$, let $v_n$ be the infimum of $\int f_1\,d\pi_1+\int f_2\,d\pi_2$
over Borel probability measures $\pi_1$ on $[0,1]\times[-1,1]$ and $\pi_2$
on $[-1,1]\times[0,1]$ whose $y$-marginals have equal moments of degrees
$0,\dots,n$. This is the \emph{local-measure relaxation}: each bag is an
actual measure, and only the separator moments are shared.

\begin{proposition}[Local-measure value]\label{prop:moment-matching}
For every $n\ge0$, $v_n=-2E_n(\varphi)$, and the infimum is attained.
Moreover, $\rhomod_r\le\rhopre_r\le v_{2r}$ for every $r\ge1$.
\end{proposition}

\begin{proof}
Let $\mu_1,\mu_2$ be the $y$-marginals of a feasible pair. By
\eqref{eq:shp-fibers}, $\int f_1\,d\pi_1+\int f_2\,d\pi_2\ge
-\bigl(\int\varphi\,d\mu_1-\int\varphi\,d\mu_2\bigr)\ge-2E_n(\varphi)$ by
\cref{lem:moment-matching-general}. Conversely, let $(\mu_1,\mu_2)$ attain
the maximum in \cref{lem:moment-matching-general}, let $\pi_1$ be the law
of $(\max\{Y,0\},Y)$ for $Y\sim\mu_1$, and let $\pi_2$ be the law of
$(Y,\max\{-Y,0\})$ for $Y\sim\mu_2$. By \eqref{eq:shp-fibers} this pair has
objective $-2E_n(\varphi)$.

For the second claim, integration against a feasible pair with $n=2r$ is a
feasible preordering family of order $r$: each local functional is
integration against a measure on its rectangle, hence nonnegative on every
square times a product of generators, and separator moments agree through
degree $2r$. \Cref{lem:weak-duality} gives $\rhomod_r\le\rhopre_r$.
\end{proof}

\subsection{An explicit approximation lower bound}\label{sec:shp-fejer}

\begin{lemma}[Lower bound for $\varphi$]\label{lem:fejer-lower}
For every $n\ge0$ there are Borel probability measures $\mu_+,\mu_-$ on
$[-1,1]$ with equal moments of degrees $0,\dots,n$ and
\begin{equation}\label{eq:shp-fejer}
  \int\varphi\,d\mu_+-\int\varphi\,d\mu_-\ge\frac2{27\pi(n+2)^2}.
\end{equation}
Consequently $E_n(\varphi)\ge1/(27\pi(n+2)^2)$.
\end{lemma}

\begin{proof}
Let $N$ be the smallest even integer with $N\ge n+1$, so $2\le N\le n+2$.
For $N<k<3N$ let $w_k=1-\abs{k-2N}/N$, and for odd $k$ let
$\epsilon_k=-\sin(k\pi/2)\in\{\pm1\}$. Define the polynomial
\[
  q_N(y)=\sum_{\substack{N<k<3N\\k\text{ odd}}}\epsilon_kw_k\,T_k(y)
\]
and the signed measure $d\lambda=q_N\,d\arcs$ on $[-1,1]$.

\emph{Annihilation.} Every $T_k$ in $q_N$ has $k\ge N+1>n$, so
$\int y^j\,d\lambda=0$ for $j\le n$ by orthogonality; in particular
$\lambda([-1,1])=0$.

\emph{Total variation.} Let $F_N(t)=\sum_{\abs j<N}(1-\abs j/N)e^{\mathrm{i}jt}
\ge0$ be the Fejér kernel, which has integral one against $dt/(2\pi)$, and
let $Q_N(\theta)=\sin(2N(\theta-\frac\pi2))F_N(\theta-\frac\pi2)$. Using
$2\sin(2Nt)\cos(jt)=\sin((2N+j)t)+\sin((2N-j)t)$,
\[
  Q_N(\theta)=\sum_{k=N+1}^{3N-1}w_k\sin\bigl(k(\theta-\tfrac\pi2)\bigr)
  =\sum_kw_k\bigl[\cos(\tfrac{k\pi}2)\sin(k\theta)-\sin(\tfrac{k\pi}2)
  \cos(k\theta)\bigr].
\]
The even part $\frac12(Q_N(\theta)+Q_N(-\theta))$ is
$\sum_{k\text{ odd}}\epsilon_kw_k\cos(k\theta)=q_N(\cos\theta)$. Since
$\cos\theta$ pushes $d\theta/(2\pi)$ on the circle forward to $\arcs$,
\[
  \int\abs{q_N}\,d\arcs=\int_{-\pi}^{\pi}\abs{q_N(\cos\theta)}
  \frac{d\theta}{2\pi}\le\int_{-\pi}^{\pi}\abs{Q_N(\theta)}\frac{d\theta}{2\pi}
  \le\int_{-\pi}^{\pi}F_N\Bigl(\theta-\frac\pi2\Bigr)\frac{d\theta}{2\pi}=1.
\]

\emph{Pairing with $\varphi$.} For odd $k$,
\[
  \int\varphi\,T_k\,d\arcs=\frac1\pi\int_0^{\pi/2}\cos^2\theta\cos(k\theta)
  \,d\theta=\frac{\sin(k\pi/2)}{2\pi}\Bigl(\frac1k-\frac{k}{k^2-4}\Bigr)
  =-\frac{2\sin(k\pi/2)}{\pi k(k^2-4)},
\]
using $\cos^2\theta=\frac12(1+\cos2\theta)$ and
$\sin((k\pm2)\pi/2)=-\sin(k\pi/2)$. Since $\sin^2(k\pi/2)=1$ for odd $k$,
\[
  \int\varphi\,d\lambda=\frac2\pi\sum_{\substack{N<k<3N\\k\text{ odd}}}
  \frac{w_k}{k(k^2-4)}.
\]
Every $k$ in the sum satisfies $3\le k<3N$, so $0<k(k^2-4)<27N^3$. The odd
$k$ are $k=2N\pm j$ with $j$ odd and $1\le j\le N-1$; since $N$ is even
there are $N/2$ such $j$, with sum $N^2/4$, so
$\sum_{k\text{ odd}}w_k=2\sum_j(1-j/N)=N/2$. Hence
$\int\varphi\,d\lambda\ge(2/\pi)(N/2)/(27N^3)=1/(27\pi N^2)$.

\emph{Measures.} Let $\lambda=\lambda_+-\lambda_-$ be the Jordan
decomposition. Since $\lambda([-1,1])=0$ and $\abs\lambda([-1,1])\le1$,
both parts have the same mass $t\in(0,\frac12]$; $t>0$ because
$\int\varphi\,d\lambda>0$. Put $\mu_\pm=\lambda_\pm/t$. These are
probability measures with equal moments through degree $n$, and
$\int\varphi\,d\mu_+-\int\varphi\,d\mu_-=t^{-1}\int\varphi\,d\lambda\ge
2/(27\pi N^2)\ge2/(27\pi(n+2)^2)$. The bound on $E_n(\varphi)$ follows from
the elementary half of \cref{lem:moment-matching-general}.
\end{proof}

\subsection{The sharp rate}\label{sec:shp-rate}

\begin{theorem}[A fixed quadratic with sparse gap of exact order $r^{-2}$]
\label{thm:quad-sharp}
For problem \eqref{eq:shp-problem}, with $\fstar=0$:
\begin{enumerate}[label=(\alph*)]
\item \emph{(Lower bound by local measures.)} For every $r\ge1$ there are
  probability measures $\pi_1$ on $[0,1]\times[-1,1]$ and $\pi_2$ on
  $[-1,1]\times[0,1]$ whose $y$-marginals have equal moments through
  degree $2r$ and
  \[
    \int f_1\,d\pi_1+\int f_2\,d\pi_2\le-\frac1{54\pi(r+1)^2}.
  \]
  Consequently
  $\rhomod_r\le\rhopre_r\le v_{2r}\le-1/(54\pi(r+1)^2)$. The same bound
  holds for every relaxation whose feasible set contains the integration
  functionals of such pairs.
\item \emph{(Upper bound, preordering.)} For every $r\ge2$,
  \[
    \rhopre_r\ge-\frac{30}{2(\lfloor r/2\rfloor+1)^2+1}\ge-\frac{60}{r^2}.
  \]
\item \emph{(Ordinary module.)} For all sufficiently large $r$,
  $\rhomod_r\ge-c\log^3r/r^2$ for an absolute constant $c$. Thus
  $-\rhomod_r$ lies between $\Omega(r^{-2})$ and $O(\log^3r/r^2)$.
\item \emph{(Dense exactness.)} With the generators $g_x,g_y,g_z$ on the
  single bag $\{x,y,z\}$, $f$ belongs to the dense truncated preordering of
  order two:
  \begin{equation}\label{eq:shp-dense}
    f=(y-x+z)^2+2\bigl((xz)^2+z^2g_x+x^2g_z+g_xg_z\bigr).
  \end{equation}
\item \emph{(No exact polynomial separator.)} There are no polynomials
  $q_1(x,y)$ and $q_2(y,z)$, nonnegative on $[0,1]\times[-1,1]$ and
  $[-1,1]\times[0,1]$ respectively, with $f=q_1+q_2$.
\end{enumerate}
In particular, $-\rhopre_r=\Theta(r^{-2})$, and the same order holds for the
local-measure relaxation: $-v_{2r}=\Theta(r^{-2})$.
\end{theorem}

\begin{proof}
(a) Apply \cref{lem:fejer-lower} with $n=2r$, lift $\mu_+$ and $\mu_-$ as
in the proof of \cref{prop:moment-matching}, and use \eqref{eq:shp-fibers}:
the objective is
$-(\int\varphi\,d\mu_+-\int\varphi\,d\mu_-)\le-2/(27\pi(2r+2)^2)=
-1/(54\pi(r+1)^2)$. The chain of inequalities is
\cref{prop:moment-matching}.

(b) In the transformed coordinates of \cref{sec:shp-example},
\[
  f_1=\tfrac38+\tfrac12\hat x+\tfrac18T_2(\hat x)-\hat xy-y,\qquad
  f_2=\tfrac78+\tfrac12T_2(y)+\tfrac18T_2(\hat z)+\tfrac12\hat z+\hat zy+y,
\]
so $\Abud(f_1)=\frac12+\frac48+2+1=4$ and
$\Abud(f_2)=\frac42+\frac48+\frac12+2+1=6$. Here $w=2$ and $d=2$, so
\cref{thm:pre} applies for $r\ge2$ with $m=\lfloor r/2\rfloor+1$ and
$\Abud=10$. Since $m>r/2$, $2m^2+1>r^2/2$.

(c) The expansions in (b) give $\Cf=\frac{21}8+\frac{25}8=\frac{23}4$, and
$w=2$, $d_\infty=2$. Apply \cref{cor:mod-rate}; with (a) this gives the
two-sided statement.

(d) Expanding verifies \eqref{eq:shp-dense}. The terms $(y-x+z)^2$ and
$(xz)^2$ are squares of polynomials of degree at most two, $z^2g_x$ and
$x^2g_z$ are squares of degree two times one generator, and $g_xg_z$ is a
product of two generators with the constant square one; every term has
total degree at most four.

(e) Suppose $f=q_1+q_2$ as stated. Then $q_1-f_1=f_2-q_2$ lies in
$\R[x,y]\cap\R[y,z]=\R[y]$, so $q_1=f_1+p(y)$ and $q_2=f_2-p(y)$ for a
polynomial $p$. By \eqref{eq:shp-fibers}, nonnegativity of $q_1$ and $q_2$
on their rectangles is equivalent to $p\ge\varphi$ and $p\le\varphi$ on
$[-1,1]$. Then $p=\varphi$ on $[-1,1]$; $p$ vanishes on $[-1,0]$, hence
identically, but $\varphi(1)=1$, a contradiction.
\end{proof}

Part (e) explains part (a). A sparse certificate of $f-\fstar$ would be
such a decomposition, so none exists at any order, and a certificate of
$f+\epsilon$ requires a polynomial $p$ of degree at most $2r$ with
$\varphi\le p\le\varphi+\epsilon$ on $[-1,1]$; by \cref{lem:fejer-lower},
$\epsilon\ge2E_{2r}(\varphi)$, with this lower bound of order $r^{-2}$.
This is a property of
the hierarchy, not of the problem: the problem is solved by
\eqref{eq:shp-fibers}, and the dense certificate \eqref{eq:shp-dense} uses
the product $xz$, which crosses the two bags. The ordinary-module upper
bound in (c) carries the logarithmic factor of \cref{cor:mod-rate}; we do
not claim that the ordinary-module gap is of exact order $r^{-2}$.

With the linear generators $x,1-x,z,1-z$ in place of $g_x,g_z$, the
identity $f=(y-x+z)^2+2xz$ is already a dense certificate of degree two.
The choice of generators does not affect part (a), whose witnesses are
actual measures.

\subsection{The first two orders}\label{sec:shp-orders}

For $r=1$ and $r=2$ the sparse values can be computed exactly. Write
$\sqrt3$ explicitly and set
\begin{equation}\label{eq:shp-constants}
  a=3\sqrt3-5,\quad b=\sqrt3-1,\quad
  c_\star=\frac{3+2\sqrt3}{18},\quad
  E_\star=\frac{7\sqrt3}9-\frac43,\quad
  p_\star(y)=c_\star(y+1)(y+a)^2,
\end{equation}
so that $0<a<b<1$ and $E_\star>0$ (numerically $a\approx0.196$,
$b\approx0.732$, $E_\star\approx0.0138$).

\begin{proposition}[Exact values at orders one and two]
\label{prop:exact-orders}
For problem \eqref{eq:shp-problem}:
\begin{enumerate}[label=(\alph*)]
\item $\rhomod_1=\rhopre_1=-\frac14$, while
  $v_2=-(3-2\sqrt2)\approx-0.1716$.
\item $\rhomod_2=\rhopre_2=v_4=-2E_\star=\frac83-\frac{14\sqrt3}9\approx
  -0.027635$, and $E_3(\varphi)=E_4(\varphi)=E_\star$.
\end{enumerate}
\end{proposition}

At order one the finite SDP has a larger gap than the local-measure
relaxation: $\rhopre_1<v_2$. At order two the entire gap is caused by the
separator truncation, even for the ordinary module.

\begin{proof}
(a) At order one the product $g_xg_y$ has degree four and is absent, so the
two cones coincide and $\rhomod_1=\rhopre_1$. The identity
\begin{equation}\label{eq:shp-order-one}
  x^2-2xy+\tfrac12y^2+\tfrac12y+\tfrac18
  =\tfrac12\bigl(2x-y-\tfrac12\bigr)^2+g_x
\end{equation}
exhibits the left side in the order-one module of $B_1$. Substituting
$(x,y)\mapsto(z,-y)$ gives
$z^2+2yz+\frac12y^2-\frac12y+\frac18$ in the order-one module of $B_2$, and
the two left sides add to $f+\frac14$. Hence $\lammod_1\ge-\frac14$, and
$\rhomod_1\ge-\frac14$ by \cref{lem:weak-duality}.

For the reverse inequality, let $L_1$ be integration against the atomic
measure $\frac34\delta_{(0,-1/2)}+\frac14\delta_{(1,3/2)}$ in the variables
$(x,y)$, restricted to polynomials of degree at most two. Its moment matrix
is PSD because $L_1$ is integration against a positive measure, and
\[
  L_1(y)=0,\quad L_1(y^2)=\tfrac34,\quad L_1(g_x)=0,\quad L_1(g_y)=\tfrac14,
  \quad L_1(f_1)=-\tfrac12 .
\]
The generator tests of order one are $L_1(g_x)\ge0$ and $L_1(g_y)\ge0$, so
$L_1$ is feasible for $B_1$. The atom $(1,\frac32)$ lies outside the
rectangle. Moreover, $L_1(xy)=3/8>L_1(x)=1/4$, whereas $xy\le x$
on the rectangle. Thus no probability measure supported on the bag
represents these truncated moments. Let $L_2$ be integration against the image measure under
$(X,Y)\mapsto(y,z)=(-Y,X)$. It is feasible for $B_2$ by symmetry, its
separator moments $L_2(y)=0$, $L_2(y^2)=\frac34$ agree with those of $L_1$,
and $L_2(f_2)=\frac34\cdot\frac14+\frac14\cdot\frac14=\frac14$, the value
of $(y+z)^2=(X-Y)^2$ at the two atoms. The total objective is $-\frac14$.

For $v_2$, let $c=\sqrt2-1$ and $S_1(y)=\frac12y^2+cy$. The error
$e=\varphi-S_1$ is odd, equals $\frac12y^2-cy$ on $[0,1]$, and has extreme
values $-\frac12c^2$ at $y=c$ and $\frac12-c=\frac12c^2$ at $y=1$. Hence
$\norm e_\infty=\frac12c^2$, attained with alternating signs at
$-1,-c,c,1$. If a polynomial $q$ of degree at most two had
$\norm{\varphi-q}_\infty<\frac12c^2$, then $q-S_1=e-(\varphi-q)$ would have
the sign of $e$ at these four points, hence at least three zeros, so
$q=S_1$, a contradiction. Thus $E_2(\varphi)=\frac12c^2$ and
$v_2=-c^2=-(3-2\sqrt2)$ by \cref{prop:moment-matching}.

(b) \emph{Approximation.} Direct expansion gives the identities
\begin{equation}\label{eq:shp-p-identities}
  p_\star(y)-y^2=c_\star\bigl(y+7-4\sqrt3\bigr)(y-b)^2,\qquad
  p_\star(y)+p_\star(-y)=y^2+2E_\star .
\end{equation}
Since $7-4\sqrt3>0$, the first gives $p_\star\ge y^2=\varphi$ on $[0,1]$,
and $p_\star\ge0=\varphi$ on $[-1,0]$ by \eqref{eq:shp-constants}. With
$\varphi(y)+\varphi(-y)=y^2$, the second gives
$p_\star(y)-\varphi(y)=2E_\star-(p_\star(-y)-\varphi(-y))\le2E_\star$. So
$0\le p_\star-\varphi\le2E_\star$ on $[-1,1]$, and $p_\star-\varphi$ takes
the values $0,2E_\star,0,2E_\star,0,2E_\star$ at
$-1,-b,-a,a,b,1$ (use the factorizations and the second identity). The
cubic $p_\star-E_\star$ therefore approximates $\varphi$ with error
$E_\star$, alternating in sign at six points. A polynomial of degree at
most four with smaller error would differ from $p_\star-E_\star$ by a
nonzero polynomial with five sign changes, which is impossible. Hence
$E_4(\varphi)=E_3(\varphi)=E_\star$.

\emph{Certificate.} Let
\[
  V=\begin{pmatrix}x(x-b)\\x(y-b)\\(y+1)(y+a)-3bx\end{pmatrix},\quad
  U=\begin{pmatrix}x-b\\y-b\end{pmatrix},\quad
  Z=(b+a)x-b(y+a),
\]
\[
  G_0=\begin{pmatrix}
  2-\sqrt3&\frac{3\sqrt3-7}4&\frac{\sqrt3-1}{12}\\[2pt]
  \frac{3\sqrt3-7}4&\frac76&-\frac{\sqrt3+1}{12}\\[2pt]
  \frac{\sqrt3-1}{12}&-\frac{\sqrt3+1}{12}&\frac1{12}+\frac{\sqrt3}{18}
  \end{pmatrix},\quad
  G_1=\begin{pmatrix}2-\sqrt3&\frac{3\sqrt3-7}4\\[2pt]
  \frac{3\sqrt3-7}4&1\end{pmatrix},\quad
  k_\star=\frac16+\frac{7\sqrt3}{72}.
\]
(The third entry of $V$ is $(y+1)(y+a)-\frac{(b+1)(b+a)}bx$, simplified
using $(b+1)(b+a)/b=3b$.) Direct expansion gives the polynomial identity
\begin{equation}\label{eq:shp-gram}
  x^2-2xy+p_\star(y)=V^{\top}G_0V+g_x\,U^{\top}G_1U+g_y\,k_\star Z^2 .
\end{equation}
The leading principal minors of $G_0$ are $2-\sqrt3$,
$(35\sqrt3-58)/24$ and $(38-15\sqrt3)/864$; those of $G_1$ are $2-\sqrt3$
and $(13\sqrt3-22)/8$. All are positive, since $\sqrt3<2$,
$35^2\cdot3=3675>58^2=3364$, $15\sqrt3<30<38$, and
$13^2\cdot3=507>22^2=484$; also $k_\star>0$. Hence $G_0,G_1$ are positive
definite, their quadratic forms are sums of squares, and every term of
\eqref{eq:shp-gram} has total degree at most four with at most one
generator. So the left side of \eqref{eq:shp-gram} is in the order-two
ordinary module of $B_1$. Substituting $(x,y)\mapsto(z,-y)$, which fixes
$g_y$, puts $z^2+2yz+p_\star(-y)$ in the order-two module of $B_2$. By
\eqref{eq:shp-p-identities} the two left sides add to $f+2E_\star$, so
$\lammod_2\ge-2E_\star$ and $\rhomod_2\ge-2E_\star$.

\emph{Witness.} Let $\mu_1$ put masses
\[
  w_{-1}=\frac{8-4\sqrt3}9,\qquad w_{-a}=\frac5{18}+\frac{\sqrt3}6,\qquad
  w_b=\frac{5\sqrt3}{18}-\frac16
\]
on the points $-1,-a,b$. Direct substitution of \eqref{eq:shp-constants}
gives
\[
  w_{-1}+w_{-a}+w_b=1,\qquad
  -w_{-1}-a w_{-a}+b w_b=0,\qquad
  -w_{-1}-a^3 w_{-a}+b^3 w_b=0.
\]
The masses are positive because $\sqrt3<2$ and $\sqrt3>1$. Thus the
first and third moments of $\mu_1$ vanish. Let $\mu_2$ be the reflection of
$\mu_1$ under $y\mapsto-y$; then $\mu_1$ and $\mu_2$ have equal moments
through degree four. Lifting as in \cref{prop:moment-matching} gives local
measures on the two rectangles with objective
$-\int\varphi\,d\mu_1+\int\varphi\,d\mu_2=-w_bb^2+w_{-1}+w_{-a}a^2
=-2E_\star$. Hence $v_4\le-2E_\star$, and
\[
  -2E_\star\le\rhomod_2\le\rhopre_2\le v_4\le-2E_\star
\]
by \cref{prop:moment-matching}. This also gives $v_4=-2E_4(\varphi)$
directly.
\end{proof}

The proof of \cref{prop:exact-orders} uses neither numerical SDP solutions
nor strong SDP duality: certificates bound the values from below and explicit
functionals bound them from above. The values for $r=1,2$ agree, for both
cones, with $-2E_{2r}(\varphi)$ at $r=2$ and exceed it in magnitude at
$r=1$.

\begin{remark}[An optimal majorant need not be certifiable at the same order]
\label{rem:shp-majorant}
The quadratic $\frac12(y+c)^2=S_1+\frac12c^2$, with $c=\sqrt2-1$, is a
majorant of $\varphi$ of minimal uniform width among polynomials of degree
at most two, so $q(x,y)=x^2-2xy+\frac12(y+c)^2$ is nonnegative on
$[0,1]\times[-1,1]$. It does not belong to the order-one module of $B_1$.
Suppose $q=\sigma_0+\gamma_xg_x+\gamma_yg_y$ with $\sigma_0$ a sum of
squares of affine polynomials and $\gamma_x,\gamma_y\ge0$. At $(0,-c)$ we
have $q=0$, $g_x=0$ and $g_y>0$, so $\gamma_y=0$ and every affine factor
of $\sigma_0$ vanishes there; write them as $\alpha_ix+\beta_i(y+c)$.
Comparing the coefficients of $y^2$, $xy$, $x$ and $x^2$ gives
$\sum\beta_i^2=\frac12$, $\sum\alpha_i\beta_i=-1$, $\gamma_x=2c$ and
$\sum\alpha_i^2=1+2c$. The Cauchy--Schwarz inequality then requires
$1\le\frac12(1+2c)=\sqrt2-\frac12<1$, a contradiction. Thus polynomial
nonnegativity of a local term does not imply membership in the local cone
of the same degree, and \cref{prop:moment-matching} cannot be transferred
to the SOS hierarchy at a fixed order without a degree-specific
certificate such as \eqref{eq:shp-gram}.
\end{remark}

Whether $\rhomod_r=\rhopre_r=-2E_{2r}(\varphi)$ for every $r\ge2$ is open;
floating-point SDP solutions for $3\le r\le6$ are numerically close to
$-2E_{2r}(\varphi)$, but they are not certified and we do not use them.

\subsection{Relation to prior work}\label{sec:shp-prior}

\begin{remark}[What is prior in \cref{thm:quad-sharp}]\label{rem:shp-prior}
Failure of finite exactness of sparse hierarchies, explained by separator
functions that are not polynomials, is known. Nie, Qu, Tang and
Zhang~\cite{nie2026-sparse-tightness}, Example~6.7, give a quartic two-bag
problem whose dense hierarchy is exact at order two and whose sparse
hierarchy is never exact; its separator function is rational. Grimm,
Netzer and Schweighofer~\cite{grimm2007-structured-sparsity} obtain sparse
representations under strict positivity by approximating separator fiber
minima by polynomials, which the strict margin permits. The identity of
\cref{lem:moment-matching-general} is known~\cite{han2018-local-moment-matching}.
Baldi and Slot~\cite{baldi2024-putinar-hypercube} prove a degree lower bound
for the dense ordinary module on a polynomial that lies in the dense
preordering at degree four; that obstruction is of a different kind,
coming from the missing generator products rather than from restricted
bags, and neither lower bound implies the other.

What \cref{thm:quad-sharp} adds, to our knowledge, is a quantitative
version: a fixed quadratic objective with a piecewise quadratic separator
function, a matching upper bound for the sparse preordering, and explicit
local-measure witnesses showing that the rate $r^{-2}$ is optimal for that
hierarchy even when every local functional is a measure.
\cref{prop:exact-orders} is a small-order benchmark for this instance, not
a general theorem.
\end{remark}
